\section{Polynomial transport and the invertible coefficient map}
\label{sec:heading-history-handoff}
The primary comparison of Section~\ref{sec:primary-calibration-producer}
uses a particular coefficient map. We prove the two-sided bound needed
to iterate it, and then combine the primary, regular and forced
insertions on one transition. The coefficient norm throughout is the
sum of the absolute values of the coordinates.

\subsection{Exact polynomial transport on the cover}
Fix the positive blocking penalty $k>0$. Let $L$ be the block side,
$u=L^{-1}$, and let $j,t\ge4$ be the
adjacent fine and coarse kernel levels. Write $A_x^{ab}(z)$ for the
real one-form transport row at the fine site $x$, with output direction
$a$, input direction $b$, and coarse displacement $z$. The actual
input site in this row is $\lfloor x/L\rfloor+z$; the floor here
means the centered block quotient. The strip estimate gives constants
$\sigma,C>0$, selected before $L$, such that
\[
 \sum_z |A_0^{ab}(z)|e^{\sigma|z|_1}\le Cu.
\]
The analogous rows at other residues are exponentially summable too.
We first show that, for every vector-valued polynomial $p$ of degree
less than $j$,
\begin{equation}\label{eq:heading-polynomial-row}
 (Ap)(x)_a=\sum_b\sum_z A_0^{ab}(z)p_b(ux+z).
\end{equation}
This is an exact identity, including for $x$ outside the central block.

Here are the low-frequency details. Denote the analytic alias column
of Section~\ref{sec:kernel-transport-identities} by $\mathcal A_l(q)$,
in its rescaled one-form convention. Fourier reconstruction and
$x=L\lfloor x/L\rfloor+\operatorname{rem}(x)$ give
\[
 e^{iq\cdot\lfloor x/L\rfloor}\widehat A_x^{ab}(q)
 =u\sum_l e^{i(q+2\pi l)\cdot x/L}\mathcal A_l(q)_{ab}.
\]
For a nonzero alias $l$, both terms in the regular one-form formula
have the factor
\[
 z_l(q)^j,\qquad
 z_l(q)=\operatorname{area}_L(-(q+2\pi l)/L),\qquad z_l(0)=0.
\]
Indeed the transposed filter contains the $j$th power of this area
factor, and the nonprincipal gain contains the same factor. Removing
it leaves the analytic port, rail and gain factors from the component
estimates. Their denominators are nonzero at the origin. Thus
$\mathcal A_l=z_l^j F_l$ with $F_l$ analytic there. Every derivative
of order less than $j$ of this term, even after multiplication by
the displayed exponential, vanishes at zero: in the product of $j$
vanishing factors some factor remains undifferentiated. The principal
column has $\mathcal A_0(0)=I$, obtained by the regular principal
formula at zero. In particular
\begin{equation}\label{eq:heading-row-mass}
                       \sum_z A_0^{ab}(z)=u\delta_{ab}.
\end{equation}

Restrict $q$ to $s\xi$ for a real covector $\xi$, and differentiate
$r<j$ times in $s$. Exponential summability justifies differentiation
of the Fourier series: its differentiated terms are bounded by a
polynomial in $|z|_1$ times the summable exponential. All nonprincipal
aliases disappear. Comparing the resulting principal expression at
$x$ and at zero, after inserting $e^{is\xi(ux)}$ at zero, gives
\[
 \sum_z A_x^{ab}(z)\xi(\lfloor x/L\rfloor+z)^r
       =\sum_z A_0^{ab}(z)\xi(ux+z)^r.
\]
Polarization of the symmetric $r$-linear forms gives the same identity
for products of $r$ possibly different covectors. Such products span
the homogeneous polynomials. Summing the homogeneous parts and the
finitely many colour coordinates proves
\eqref{eq:heading-polynomial-row}.

Let $p$ now have homogeneous degree $r<j$. Expanding $p(ux+z)$
multilinearly, a term containing $s$ copies of $x$ has a factor $u^s$.
Its remaining $r-s$ factors are bounded by a constant times
$(1+|z|_1)^{r-s}$. Consequently its coefficient tensor has norm at most
\[
 C_r\|p\|u^{s+1},\qquad 0\le s\le r.
\]
The constants are independent of $L$: use
$\sup_{v\ge0}(1+v)^{r-s}e^{-\sigma v}<\infty$ and the row bound.
Passing from diagonal polynomial values to symmetric tensor norms
costs only a fixed finite-dimensional constant. This follows by
compactness of the tensor unit sphere and injectivity of the diagonal
map, or directly by polarization. The degree-$r$ term is exactly
$u^{r+1}p(x)$ by \eqref{eq:heading-row-mass}. Adding homogeneous
parts proves that, if $\deg p<a\le j$, then
\begin{equation}\label{eq:heading-polynomial-leading}
 \deg(Ap)<a,\qquad \deg(Ap-u^a p)<a-1,
\end{equation}
and the degree-$s$ output tensor is bounded by
$C_a u^{s+1}$ times the sum of the input tensor norms.
Lower-degree input terms do not affect the asserted highest-degree
identity.

\subsection{The two gradings of the primary matrix}
Let $e_i$ be the canonical coordinate basis of the degree-$d$ frame,
with $e_0$ the unit, and write $(g_i,\ell_i)$ for spatial grade and
field-amplitude grade. Both lie between zero and $d$. Let $R$ be the
unscaled polynomial coefficient transport used in
Section~\ref{sec:primary-calibration-producer}. We claim
\begin{align}
 R_{bi}&=u^{g_i}\delta_{bi} &&(g_b\le g_i),
                          \label{eq:heading-triangular-receipt}\\
 R_{bi}&=0 &&(\ell_b<\ell_i),\qquad
 |R_{bi}|\le D u^{g_i}.    \label{eq:heading-receipt-price}
\end{align}
Here $D\ge1$ is chosen before $L$, for the fixed frame and kernel
levels, and we assume $d\le j$.

To verify the first identity, differentiate the local holonomy
display in directions of engineering weights $w_1,\ldots,w_n$.
A block of a chain-rule partition has weight
$a=\sum_{r\text{ in block}}w_r$. Its unsmoothed potential is
polynomial of degree less than $a$. By
\eqref{eq:heading-polynomial-leading}, the smoothed potential is
$u^a$ times that potential plus a polynomial of degree less than
$a-1$. Expand the outer multilinear derivative of the $i$th section
in these sums. If the total weight is at most $g_i$, every term using
one lower-degree correction has weight strictly below $g_i$ and
vanishes by the primary section's lower-grade cancellation: its
multilinear derivative is zero on polynomial tuples of total
engineering weight less than $g_i$. Indeed its local-section tests
read the homogeneous label $e_i$, so all smaller spatial grades
vanish (Section~\ref{sec:local-sections}); the filtered
gauge-coordinate argument of
Section~\ref{sec:polynomial-moment-attenuation} then gives exactly
this lower-weight assertion. No vanishing of its grade-$g_i$ test
is claimed. The remaining term
has factor $u^{\sum_r w_r}$. The unsmoothed section reads exactly
$e_i$, by the defining local-section identity. Extracting its
homogeneous part of output bidegree $(g_b,\ell_b)$ proves
\eqref{eq:heading-triangular-receipt}. This also covers order zero.

For the second identity, the section has ordinary derivatives zero
through order $\ell_i-1$ at the flat field. In a derivative of total
order less than $\ell_i$ of its composition, the outer derivative
has no more arguments than that total order. It therefore vanishes.
This proves the amplitude-grade restriction, regardless of spatial
grade.

For the uniform bound, use the homogeneous decomposition following
\eqref{eq:heading-polynomial-leading} in every inner block of the
chain rule. A term of polynomial degree $s$ in that block has weight
$s+1$ and coefficient norm at most $C u^{s+1}$. Expand the outer
multilinear derivative. Terms whose weights sum to less than $g_i$
vanish; every other term has a factor at most $u^{g_i}$, since
$0<u\le1$. The fixed stencil, bounded derivative order and finite
partition set give a constant for each polynomial test. A finite
family of these tests controls the coefficient norm, as proved in
Section~\ref{sec:local-splitting}. Taking the maximum over the
finitely many columns and tests proves the last inequality in
\eqref{eq:heading-receipt-price}.

Let $f,c\in(0,1]$ be the fine and coarse small-field thresholds,
with $c/f\le\pi$ and $f/c\le\pi$, where $\pi\ge1$. The actual
coefficient map $B$ has entries
\[
                      B_{bi}=c^{\ell_b}f^{-\ell_i}L^{g_i}R_{bi}.
\]
If $\ell_b\ge\ell_i$, then
$c^{\ell_b}/f^{\ell_i}\le(c/f)^{\ell_i}\le\pi^d$; otherwise
the entry is zero. Set $a=\pi^dD$. We have
\[
 a\ge1,\qquad |B_{bi}|\le a,\qquad
 |B_{ii}|=(c/f)^{\ell_i}\ge a^{-1},
\]
and an off-diagonal nonzero entry requires $g_i<g_b$.
The unit is fixed. The actual polynomial transport commutes with
signed coordinate motions and conjugation: change variables in its
absolutely convergent real kernel sum, use the covariance of the
local section, and then take its finite jet. The three diagonal
scalings in $B$ preserve those actions because the bidegrees do.

\subsection{A bound for the inverse without a perturbative assumption}
Let $n$ be the number of coordinates and put $K=a+na^2\ge1$.
For $y=Bv$, induction on the grade proves
\[
                |v_i|\le K^r\|y\|_1\quad(g_i<r).
\]
The claim is empty for $r=0$. If $g_i<r+1$, the other nonzero
entries in row $i$ come from grades less than $g_i$, hence less
than $r$. Isolating the diagonal entry and applying the induction
hypothesis gives
\[
 |v_i|\le a\left(|y_i|+\sum_{b\ne i}|B_{ib}v_b|\right)
 \le a\|y\|_1+na^2K^r\|y\|_1
 \le K^{r+1}\|y\|_1.
\]
Summing with $r=d+1$, and bounding the forward map row by row, yields
\begin{equation}\label{eq:heading-both-prices}
 \|Bv\|_1\le C_B\|v\|_1,
 \qquad \|v\|_1\le C_B\|Bv\|_1,
 \qquad C_B=n(a+na^2)^{d+1}.
\end{equation}
Thus $B$ is injective and, in this finite-dimensional space,
bijective. Substituting $v=B^{-1}w$ gives the same bound for its
inverse. The unit column and grade-zero row show also that $B$
commutes with deletion of the unit coordinate, denoted by
$\operatorname{core}$. In particular this argument requires no
smallness of the off-diagonal entries.

\section{Three elementary insertions on one history transition}
\label{sec:three-port-history}
Fix one funded interpolating transition, its point on the interpolation
interval, and its fine and coarse representations $R_f,R_c$.
Retain the actual local slots and their support and inventory
identifications. The marked derivative fixes one map $\mathcal D$
from real elementary insertions to coarse forced arrays, as in
Section~\ref{sec:projection-payment-producer}. The equalities
identifying $R_f,R_c$ with adjacent entries of the chosen history
are part of these data. None of the selections below changes that
transition or its input and output representations.

Write $\mathcal T$ for its density transfer, $v=L^d$ for the
physical factor, $\eta=L^{-2/5}$ for its small error allowance,
and $\kappa_f,\kappa_c$ for the square roots of the fine and coarse
forced allowances. They satisfy $0<\kappa_f,\kappa_c\le1$.
For a support $X$ and majorant $h\ge0$, its fine mass is
\[
 m_f(X,h)=e^{b_f C_{B_f}(X)}h.
\]
The frame's finite motion price is $C_{\rm orb}\ge1$. Let $C$
be the selected coefficient allowance, at least the price $C_B$
in \eqref{eq:heading-both-prices} and the graded remainder price
in the primary estimate. All these choices precede the final
coupling threshold.

\subsection{Keeping the three exact identities}
For an inserted density $F$, the local splitting constructs a
coefficient field $Q$ and a response $\psi$ such that
\begin{equation}\label{eq:history-three-identities}
 \begin{split}
 \operatorname{dens}(Q)&=v\mathcal T F,\\
 \alpha&=\sum_r z_r\operatorname{ord}(M_r),\qquad
 \omega=\sum_r z_r\operatorname{flag}(M_r),\\
 \mathcal P_c(Q_{\rm forced})-\operatorname{hook}(\psi)
       &=v\sum_r z_r\mathcal P_c(\mathcal D M_r).
 \end{split}
\end{equation}
Here $\alpha$ is the input exponential entry, $\omega$ its input
forced profile, $M_r$ are real elementary insertions, and
$\mathcal P_c$ groups output coefficients with the same support and
inventory. The ordinary and flagged entries are separate functions;
the density combines them with the background density and gas.
The first identity follows from the actual integral expansion and
the exact subtraction identity of Section~\ref{sec:local-splitting}.
The last is the retained singleton identity
\eqref{eq:projection-profile-retained}. The middle identities simply
specify the input decomposition; they are not inferred from the
first.

For a complex elementary insertion use its real and imaginary
parts with the same majorant, and coefficients $1,i$. Their input
cost is at most twice their common mass. The response $\psi$ is
supported on sets containing the origin, is local, gauge invariant
and measurable, and vanishes on healthy configurations of its
support. If its reserved mass bound is $a_0$, the local splitting
supplies remainder budgets $a_1,a_2\ge0$ with
\[
 \operatorname{rem}_{\kappa_c}(Q)\le a_1+a_2,
 \qquad C_{\rm turn}a_1+a_2\le r,
 \qquad a_0\le a_1\le r,
\]
where $C_{\rm turn}\ge1$. More precisely the first two budgets
retain the separate ordinary and forced estimates, so the same
bound $r$ holds after the allowed finite motions. These assertions
are the summation and support-join bounds already proved in the
local splitting; the response is the sum of its rebates. Absolute
majorants justify both its grouping by support and the profile
identity in \eqref{eq:history-three-identities}.

\subsection{The three choices and their prices}
For the primary insertion with coordinate $e_i$, $i\ne0$, take
its local section on the canonical stencil. Its exponential entry
is its masked section, its input forced profile is zero, and its
density is that entry times the fine background density. The
small-error and graded remainder estimates of
Section~\ref{sec:primary-calibration-producer}, followed by the
same local splitting, give
\[
 C_{\rm orb}\|Q.x-BL^{d-g}e_i\|_1
       \le\eta\|\operatorname{core}e_i\|_1,
 \qquad r=C\|L^{d-g}\operatorname{core}e_i\|_1.
\]
The diagonal notation $L^{d-g}$ acts coordinatewise. Let $\tau_f$
be the common fine local-section derivative bound. Its stencil is
contained in one $B_f$-covering set once the selected geometric
floor is met. Thus its input mass is at most $e^{b_f}\tau_f$.
Both real components therefore fit the common input allowance
\begin{equation}\label{eq:history-common-waybill}
                            t_f=2e^{b_f}\tau_f.
\end{equation}
This allowance is independent of the coordinate.

For a regular insertion $f$ on $X\ni0$, assume the local regularity,
gauge invariance and vanishing holonomy tests through degree $d$,
with the prescribed derivative and spatial-weight bound $h$.
Its input exponential entry is $\chi_X f$, its forced profile is
zero, and its density is $\rho_f\chi_X f$. Put $m=m_f(X,h)$.
Section~\ref{sec:regular-principal-producer} pays its principal
term; the nonprincipal and forced payments are those of
Section~\ref{sec:projection-payment-producer}. The split field
therefore satisfies
\[
          C_{\rm orb}\|Q.x\|_1\le\eta m,
          \qquad r=\eta m,
          \qquad\text{input allowance }2m.
\]
All three equalities in \eqref{eq:history-three-identities} use
this same masked insertion.

For an allowed forced slot $s$, let $f\ne0$ be its local,
gauge-invariant measurable insertion, with the slot's support
and eligibility conditions, $|f|\le h$, and origin in its support.
Put $m=m_f(X_s,h)$. Its exponential entry is zero. Its input
profile equals $f$ at the original support--inventory address
and zero at every other address. It has no ordinary principal
cluster: the marked slot is flagged, whereas that cluster requires
an ordinary insertion. The nonprincipal and forced prices thus
give the zero-leading-coefficient bound without a principal bill.
The fine forced norm charges this insertion by $m/\kappa_f$.
Since $0<\kappa_f\le1$, enlarging the already paid bound
$\eta m$ gives the required form
\begin{equation}\label{eq:history-forced-price}
 C_{\rm orb}\|Q.x\|_1\le\eta m/\kappa_f,
 \qquad r=\eta m/\kappa_f,
 \qquad\text{input allowance }2m.
\end{equation}
The input allowance remains $2m$; it is not multiplied by
$\kappa_f^{-1}$. The output remainder uses $\kappa_c$ through
the local splitting. The two square roots serve different roles.
For $f=0$ take zero fields and zero response, so no support
claim for a nonzero insertion is needed.

\subsection{One finite set of insertion types for every grade}
On the fixed finite fine torus, form a set containing the ordinary
slot selected for every support containing the origin. Adjoin the
forced slots selected by the retained slot realization for all
support--inventory tags. Both indexing sets are finite, so their
union is finite even if different indices select the same slot.
This set contains all three input types above. It depends on the
representation and its slot realization, not on the coefficient
vector later to be transported.

The map $\mathcal D$ was fixed before choosing the degree $d$.
For each $d$ with $d+3$ at most the coarse kernel level, apply the
three estimates and \eqref{eq:heading-both-prices} with its own
frame. Its common allowance is \eqref{eq:history-common-waybill},
using the stage's local-section price, which bounds all those
grades. The exact history identifications replace the interpolating
input and output by $R_{N,m+1}$ and $R_{N,m}$ in every equality.
This is substitution in fixed maps; it selects no new transition.
The transition remains sound because its unmarked density identity,
singleton derivative and motion identities were established for
this interpolating point.

Finally the common scale payments give, with $s_f$ the fine column
allowance and $K_f$ the next coefficient budget,
\[
 \eta(1+s_f)C_{\rm orb}\le L^{-3/8},\qquad
 (C+\eta s_f)C_{\rm orb}\le K_f.
\]
These are exactly the two purchase inequalities in
\eqref{eq:ex-interpolation-buy}. Together with the selected bank
and room bounds there, they allow
Proposition~\ref{prop:ex-selected-interpolation} to combine all
columns into one linear coefficient family. At a fixed stage there
are only finitely many permitted earlier stages and grades. Taking
the maximum of their coupling thresholds supplies all these
estimates simultaneously, and any subsequent increase of those
thresholds preserves them.

This proves the three elementary estimates, the inverse-map bound,
and their common-history use relative to the retained funded
interpolation, its slot realization and the common scale payments.
The construction of those nonlinear interpolation data remains the
separate path-construction problem.
